\scanpage{72}
% archive sheet number: 71
$J_{\sigma,\sigma'}$, \deleted{\unsure{on a si}} les $a\in A(\gamma)$ et les $a'\in A(\gamma')$ sont distincts. \deleted{\unsure{Le}} le produit des $Y_a,Y_{a'}$ correspondants est aussi le monôme correspondant à une structure \deleted{polygonale} \inserted{de contour} sur $J(\sigma,\sigma')$ (dont les arêtes sont les ens de la forme $\{\unsure{i|x},\sigma i\}$ et $\{i,\sigma'i\}$), dont les composantes connexes sont d'ordre pair $\geq4$. Le couple $(\sigma,\sigma')$ est déterminé, à symétrie près \inserted{$(\sigma,\sigma')\longmapsto(\sigma',\sigma)$ sur chaque composante conn. de $J(\sigma,\sigma')$}, par la structure de graphe correspondante -- et on trouve \inserted{ainsi} toutes les structures de graphe $\gamma\in\Gamma$ pour lesquelles \deleted{il n'y a} $\unsure{\nu|v}_1(\gamma)=0$, et \deleted{les cycles} $\nu_*^+(\gamma)=\nu_*(\gamma)$. Un $=$ monôme intervient $2^{\nu_*(\gamma)}=2^{\nu_*^+(\gamma)}$ fois.

Donc on trouve
\[
\begin{aligned}
B^2-(-1)^m\Delta &\equiv
\sum_{\substack{\gamma\in\Gamma\\\boxed{\substack{\nu_1(\gamma)=0\\\nu_*(\gamma)=\nu_*^+(\gamma)>0}}}}
2^{\nu_*(\gamma)}M_\gamma\qquad\mathrm{mod}\ 4\\
&\qquad+2\sum_{\gamma\in\Gamma_1}(-1)^{n+\nu_2(\gamma)}M_\gamma\\
&\equiv2\sum_{\substack{\gamma\in\Gamma_1\\\deleted{\nu_1(\gamma)=0}\\\deleted{\nu_*(\gamma)=\nu_*^+(\gamma)=1}}}M_\gamma\qquad\not\equiv0\quad(4)!\\
&\qquad\deleted{\unsure{+}}+2\sum(-1)^{\varepsilon(\gamma)}M_\gamma\\
&=2\sum M_\gamma\bigl(\underbrace{1+(-1)^{\varepsilon(\gamma)}}_{0\ \mathrm{ou}\ 2}\bigr)\equiv0\quad(4)
\end{aligned}
\]
